\documentclass[draft]{OCR}
% \documentclass[final]{OCR}

%% -------------------------------------------------------------------
%% -------------------------------------------------------------------
%%% Try not include any additional package in this template. However...
%% -------------------------------------------------------------------
%% -------------------------------------------------------------------
%%% Siunitx: You can add this package if would like to express the
%%% units in the SI system correctly. Of course, you can always do it
%%% manually.
%\usepackage{siunitx}
%% -------------------------------------------------------------------
%% -------------------------------------------------------------------
%%% Flushend: You can add this package to automatically balance the
%%% final page, but if things go awry (e.g. section contents appearing
%%% out-of-order or entire blocks or paragraphs are coloured), remove
%%% it!
% \usepackage{flushend}

\title{Manuscript Title}
\runningtitle{OCR Template}

%%% Use the \authfn to add symbols for additional footnotes, if any. 1
%%% is reserved for correspondence emails; then continuing with 2 etc
%%% for contribtions.
\author[1]{Firstname Lastname~\orcidlink{0000-0009-2217-3853}}
\author[2 \cauthfn{1}]{Secondname Lastname}
\author[2]{Thirdname Namelast~\orcidlink{0000-0009-2217-9999}}
\author[3]{Fourthname Namelast}

%%% "Short" author for running page header
\runningauthor{Lastname et al.}

%%% Corresponding Author 
\corrauth{\cauthfn{1} second@univ.edu}

\affil[1]{University 1, Third Av. 104, Sector Z,
  Building B, 3º P, Sala 2222, Country C, ZIP 01452-997}
\affil[2]{University 2 of 6, Province VII, Country X, ZIP 99-66555/4.}
\affil[3]{The Third Institute, Somecountry, etc.}

%%% ----------------------------------------------------------------------------
%%% This block should only be set by an editor
%%% ----------------------------------------------------------------------------
%%% Paper category
\jpapercat{Manuscript Template}
%%% Is this a special issue? answer is {no} or {name of the special issue}
%%% If it is not a special issue, just say no:
% \jspecial{no}
%%% Otherwise do it like this:
\jspecial{Advances in Physical Oceanography \linebreak honoring Prof. Affonso Mascarenhas}
\jissn{ISSN 2675-2824}
\jdoi{http://dx.doi.org/10.1590/FOR PEER REVIEW}
\jvolume{9}
\jnumber{00000}
\jyear{2022}
\jeditor{Rubens Mendes Lopes}
\jaeditor{Ronaldo Balofelli}
\jsubmitted{29-October-2022}
\japproved{18-December-2022}
%%% ----------------------------------------------------------------------------

\begin{document}

\begin{frontmatter}
\maketitle
\begin{abstract}
  Abstract should be limited to up to 300 words and it should give a complete understanding of the
  manuscript, including the main objective(s) or hypothesis/hypotheses, summarized methods, key
  results and final or concluding remarks. Do not use subheadings in the abstract body (such as
  ``objective'' or ``results''). Do not cite references.
% \textbf{Conclusions}
\end{abstract}

\begin{keywords}
Keyword1, keyword 2, keyword 3 (Limit to 5 keywords representing the main content of the article).
\end{keywords}
\end{frontmatter}

\section{Using this Template}

This \LaTeX{} template is a general guideline of the complete Instructions for Authors of the Ocean
and Coastal Research (OCR) journal. The manuscript sections below are given for Original
Articles. Other manuscript formats such as Review Articles may have a more flexible structure. In
the case of Brief Communications, all sections below should appear in sequence, without subsections
and subsubsections. For details, please refer to our webpage (Instruction for Authors) or contact
the editorial office at ocr\_journal@usp.br.

\textbf{Please note that whilst this template provides a preview of the typeset manuscript for
  submission, it will not necessarily be exactly the final publication layout.}

\section{Introduction}

The introduction should be concise and place the manuscript in broader context. The text should
present the current state of the knowledge in the topic related to the research, sharing the
baseline information and citing key and relevant publications.  The introduction should outline the
significance and purpose of the work and clearly describe specific hypotheses being tested, in the
last paragraph(s). Authors are free to describe their objective(s) as one or more hypothesis to be
tested, or as a question to be answered, or even as an interesting/important environmental
feature/phenomenon to be described.

The Introduction must be straight to the point, and preferably
without subheadings. Authors should keep in mind that the introduction must be comprehensible to
scientists from other research areas unrelated to the topic of the manuscript.  Lines should be
numbered in the document, starting from the first paragraph of the Introduction, as shown in this
template.

\section{Methods}

This section should be described with sufficient detail to allow others to replicate and build on
published results. New methods and protocols should be described in detail while well-established
methods can be briefly described and appropriately cited. We strongly recommend that authors
reference the study/survey in time and space, if applicable. Authors should describe the main
characteristics of the surveyed area and/or the experimental design. The description of data
analyses is strongly recommended. The name and version of any relevant software or computer code
relevant to the study should be cited or made available. Include any pre-registration
codes. Research permits should be placed in this section whenever necessary.

\section{Results}

This section must address the main results gathered from the sampling procedures, survey design
and/or experimental work described in the Methods section.

\subsection{Subheadings}

Authors are free to use specific subheadings to better share the most important results that support
or reject their research hypothesis or that better describe the oceanographic features being treated
in the manuscript.

The information in the Results section should follow the same logical presentation provided in the
Methods section. Tables and figures can be used to summarize or illustrate the results. Additional
data (large tables, extra figures, images, videos, etc.) can be shared as Supplementary Material in
case the manuscript is accepted. Please correctly identify such material upon submission.

You can use \verb|\section{...}|, \verb|\subsection{...}| commands to add more sections and
subsections to your manuscript. Further sectional levels are provided by \verb|\subsubsection|,
\verb|\paragraph| and \verb|\subparagraph|.

\subsection{Equations}

Equations are one of the main points of using \LaTeX. Here are some nonsense examples:

Let $X_1, X_2, \ldots, X_n$ be a sequence of independent and identically distributed random
variables with $\text{E}[X_i] = \mu$ and $\text{Var}[X_i] = \sigma^2 < \infty$, and let
%

\begin{equation}\label{eq:LeBron}
S_n = \frac{X_1 + X_2 + \cdots + X_n}{n}
      = \frac{1}{n}\sum_{i}^{n} X_i
\end{equation}
%
denote their mean. Then as $n$ approaches infinity, the random variables $\sqrt{n}(S_n - \mu)$
converge in distribution to a normal $\mathcal{N}(0, \sigma^2)$.

\begin{eqnarray}
  S_n = & \frac{X_1 + X_2 + \cdots + X_n}{n}
  \frac{X_1 + X_2 + \cdots + X_n}{n}
  \frac{X_1 + X_2 + \cdots + X_n}{n} \nonumber\\
  & \frac{X_1 + X_2 + \cdots + X_n}{n}
  \frac{X_1 + X_2 + \cdots + X_n}{n}
    \frac{X_1 + X_2 + \cdots + X_n}{n}\nonumber\\
        & \hspace{10.5em}\frac{X_1 + X_2 + \cdots + X_n}{n}
\end{eqnarray}
Here, have some more math:
\[ \vec{\nabla}\int_0^\infty\vec{P_i}\underbrace{(\kappa,\sigma,\varsigma, \zeta)}_{\text{true
      dimens.}}dt = \oint_0^\infty
  e^{\textstyle-\frac{2\pi}{\lambda}}\vec{\nabla}\left(\frac{\partial^2\psi}{\partial
      t^2}\right)dr_\perp\]
This last unnumbered formula has nothing to do with Formula~\ref{eq:LeBron} by Reynaud LeBron.

\subsection{Preparing Figures and Tables}

Figures and graphics at a sufficiently high resolution (minimum 1000 pixels width/height, or a
resolution of 300 dpi) can be uploaded in a proper field in the submission platform. Make
sure figures have sufficient pixel definition before submission. Common formats are accepted;
however, PNG, JPEG and EPS are preferred.  Authors are encouraged to prepare figures in color (RGB
at 8-bits per channel).

The text within the figures and graphics must be in a font size large enough to be perfectly legible
even after a figure is reduced to fit the journal’s page format. Ocean and Coastal Research can
publish multimedia files in articles or as supplementary materials. Please contact the editorial
office for further information.

All Figures and Tables should have a self-explanatory caption. All table columns should have an
explanatory heading. Further direction on Figures and Tables are detailed below .


In \LaTeX, figures and tables can be added with the usual \verb|figure| and \verb|table| environments, e.g.:
~Figure \ref{fig:example} and Table \ref{tab:example}. Use \verb|figure*| and \verb|table*| if you
need a two-column wide figure or table, as in Figure \ref{fig:example:wide} and Table
\ref{tab:example:wide}. {\bfseries We do not recommend}, but you can force a figure or table to appear at a
specific text location with \verb|figurehere| and \verb|tablehere| as in this example (see Figure~\ref{fig:example_here}).

\begin{figure}[bt!] %% preferably at bottom or top of column
\centering
\includegraphics[width=\linewidth, draft=false]{example-image}
\caption{An example figure}\label{fig:example}
\end{figure}

\begin{figurehere}
\centering
\includegraphics[width=0.3\linewidth, draft=false]{example-image}
\captionsetup{font=figcaptionfont}
\caption{An example figurehere.}\label{fig:example_here}
\end{figurehere}

\begin{table}[bt!]
\caption{An example table.}\label{tab:example}
\begin{tabular}{l r l}
\toprule
Item & Quantity & Notes\\
\midrule
Widgets & 42 & Over-supplied\textsuperscript{*} \\
Gadgets & 13 & Under-supplied \\
\bottomrule
\end{tabular}
\begin{tablenotes}
\item This is a table note.
\item \textsuperscript{*}Another note.
\end{tablenotes}
\end{table}

\section{Discussion}

The Discussion must start with a statement of main findings, preferably in one sentence, without
repeating the results.

The authors must then indicate the strengths and weaknesses of their results, in perspective with
other studies, and address the meaning of their findings without entering into speculative or
circular reasoning. Unanswered questions and future research directions may also be mentioned.

Quote tables and figures only when it is essential to draw the reader's attention to one or more
important results. 

\subsection{Subheadings}

Authors are free to use specific and useful subheadings according to their manuscript requirements
to better discuss the results.

\section{Conclusion}

This section is not mandatory but can be added to stress the main findings and future directions of
the research, usually within one or two paragraphs. This section must have a strong link to the main
objectives, questions or hypotheses. Avoid sentences such as "in conclusion ..." or "in
summary...". Use the section to highlight the value of your research and position your findings
within a large context.

\section{Acknowledgments}

Acknowledgments must be brief, straight to the point. Funding agencies and other funding sources
must be disclosed, with their respective grant number(s) if necessary. Keep the original names and
acronyms of the native language of institutions and sponsors.

\section{Author Contribution}


{\bfseries Ocean and Coastal Research} follows the {\bfseries CRediT} criteria for authorship role
designation. All co-authors must have at least (1) actively participated in the discussion of
results, and (2) reviewed and approved the final version of the manuscript. Please select the
role(s) for each author as expressed on the CRediT website at https://casrai.org/credit/ and inform
them in this section, using author initials, followed by the respective role(s).  For example:

\begin{description}
\item[F.L.]: Conceptualization; Investigation; Writing – original draft; Writing – review \&
  editing;
\item[S.L.]: Methodology; Software; Formal Analysis; Investigation; Writing – review \&
  editing;
\item[T.N.]: Data Curation; Software; Writing initial draft. 
\item[F.N.]: Supervision; Resources; Project Administration; Funding Acquisition; Writing – review
  \& editing.
\end{description}

\subsection{Citations and References}

Ocean and Coastal Research follows the Harvard style of literature referencing. For speeding up the
reviewing and the publishing processes, we strongly suggest the authors to use the ‘Ocean and
Coastal Research’ BiB\TeX\/ style (\texttt{ocr.bst}) available for download.

It is the responsibility of the authors to double check the list of references and their quotations
in the main text. Preparing the references using appropriate software will assist in avoiding typing
mistakes, duplicated references and formatting issues. Citations and References in Supplementary
files are permitted provided that they also appear in the main text and in the reference
list. Citations in the main text must include the author's names and year of publication, following
the ‘Ocean and Coastal Research’ style. References to thesis, dissertations and reports are allowed
as long as a valid and active URL pointing to the full text is included, along with the date of
access.


Use the \verb|\cite| command for parenthetical citations, and \verb|\citeasnoun| command for text
citations (when using \verb|alpha-refs|).  This is a citation: \cite{Fan:2004} and here are two more
\cite{Cox:1972,Hear:Holm:Step:quan:2006}.  This is a citation: \cite*{Fan:2004} and here is another
one: \citeasnoun{Hear:Holm:Step:quan:2006}. The resulting bibliography is as follows (and in the real case should be moved to the endo of the \texttt{.tex} file:

\section{\LaTeX-specific examples}

\begin{quote}
This is a quote, if you ever want to use one. 
\end{quote}

\begin{itemize}
\item This is a bullet list.
\item Another point.
\item A third point.
\end{itemize}

\begin{enumerate}
\item This is a numbered list.
\item Another point.
\item A third point.
\end{enumerate}

This\footnote{This is the footnote text. } is a footnote.

\subsection{This is a 2nd level heading}

\subsubsection{This is a 3rd level heading}

Use \verb|\subsubsection| to get a 3rd level heading.

\paragraph{This is a 4th level heading}

Use \verb|\paragraph| to get a 4th level heading.

\subparagraph{This is a 5th level heading}

Use \verb|\subparagraph| to get a 5th level heading.

\begin{figure*}%[b!]  %% Add a [b!] if you prefer the wide image to be at the bottm of the page
\centering
\includegraphics[width=.7\textwidth, draft=false]{example-image}
\caption{An example of a wide figure. 
}\label{fig:example:wide}
\end{figure*}

\begin{table*}[bt!]
\caption{Example of a wide table with notes. Automobile land speed records (GR 5-10)}\label{tab:example:wide}
% Use "S" column identifier (from siunitx) to align on decimal point.
% Use "L", "R" or "C" column identifier for auto-wrapping columns with tabularx.
\begin{tabularx}{\linewidth}{L l l l r L}
\toprule
{Speed (mph)} & {Driver} & {Car} & {Engine} & {Date} & {Extra comments}\\
\midrule
407.447     & Craig Breedlove & Spirit of America          & GE J47    & 8/5/63  & (Just to demo a full-width table with auto-wrapping long lines) \\
413.199     & Tom Green       & Wingfoot Express           & WE J46    & 10/2/64  \\
434.22      & Art Arfons      & Green Monster              & GE J79    & 10/5/64  \\
468.719     & Craig Breedlove & Spirit of America          & GE J79    & 10/13/64 \\
526.277     & Craig Breedlove & Spirit of America          & GE J79    & 10/15/65 \\
536.712     & Art Arfons      & Green Monster              & GE J79    & 10/27/65 \\
555.127     & Craig Breedlove & Spirit of America, Sonic 1 & GE J79    & 11/2/65  \\
576.553     & Art Arfons      & Green Monster              & GE J79    & 11/7/65  \\
600.601     & Craig Breedlove & Spirit of America, Sonic 1 & GE J79    & 11/15/65 \\
622.407     & Gary Gabelich   & Blue Flame                 & Rocket    & 10/23/70 \\
633.468     & Richard Noble   & Thrust 2                   & RR RG 146 & 10/4/83  \\
763.035     & Andy Green      & Thrust SSC                 & RR Spey   & 10/15/97\\
\bottomrule
\end{tabularx}

\begin{tablenotes}
\item Source is from this website: \url{https://www.sedl.org/afterschool/toolkits/science/pdf/ast_sci_data_tables_sample.pdf}
\end{tablenotes}
\end{table*}

%% Our template requires BiBTEX.
%% You don't need to specify the bibliography style as it is set by the \documentclass{OCR}.
%% Specify your .bib file name here, without the extension. My file is named paper-refs.bib.
\bibliography{paper-refs}

\end{document}
